\section{Diseño de APIs REST: Principios de Fielding y OpenAPI}
\subsection{Los 6 principios de REST y su aplicación en el Sistema de Pre-Sustentaciones}

Roy Fielding定義 REST en su tesis doctoral (2000) a partir de restricciones arquitectónicas. A continuación se
analiza cada principio y su grado de cumplimiento en la API del Sistema de Gestión de Pre-Sustentaciones UTEQ.

\begin{enumerate}
    \item \textbf{Client-Server:} separación de responsabilidades. El frontend Angular 21 (cliente) y el backend
    Spring Boot 3.2.1 (servidor) son aplicaciones independientes que se comunican únicamente por HTTP en el puerto
    8080 y 4200. \textbf{Cumplido.}
    \item \textbf{Stateless:} cada petición contiene toda la información necesaria. El sistema utiliza autenticación
    \textbf{JWT stateless}: el servidor no guarda sesiones (\texttt{SessionCreationPolicy.STATELESS}) y cada
    petición incluye el token en la cabecera \texttt{Authorization: Bearer} o en cookie HTTP-Only.
    \textbf{Cumplido.}
    \item \textbf{Cacheable:} las respuestas pueden declarar si son cacheadas. Redis 7 está incluido en el
    \texttt{docker-compose.yml}, pero el backend no ha habilitado `@Cacheable` en los servicios REST.
    \textbf{Falta por implementar.}
    \item \textbf{Uniform Interface:} URIs de recursos y verbos HTTP semánticos. La API usa sustantivos para los
    recursos (\texttt{/api/solicitudes}, \texttt{/api/anteproyectos}, \texttt{/api/cronogramas}, \texttt{/api/evaluaciones},
    \texttt{/api/jurados}, \texttt{/api/actas}) y verbos HTTP estándar (GET, POST, PUT, DELETE).
    \textbf{Cumplido.}
    \item \textbf{Layered System:} el cliente no conoce detalles internos. La arquitectura de capas (Controller →
    Service → Repository) y la posibilidad futura de colocar un proxy/nginx delante del backend permiten cumplir
    este principio. \textbf{Cumplido parcialmente} (todavía no hay nginx en producción).
    \item \textbf{Code-On-Demand (opcional):} el servidor puede enviar código ejecutable. No se aplica en el sistema,
    lo cual es una restricción opcional del estilo. \textbf{No aplicado.}
\end{enumerate}

\subsection{Diseño de URIs REST semánticas}

Las URIs del PFC siguen la regla \textit{recursos como sustantivos y jerarquía por identificación}. Por ejemplo,
\texttt{GET /api/solicitudes/{id}} identifica una solicitud concreta, mientras que
\texttt{GET /api/solicitudes/mis-solicitudes} obtiene las solicitudes del usuario logueado extrayendo el ID desde
el contexto del token JWT sin exponer IDs en la URL. El diseño completo de recursos principales se muestra en la
Tabla~\ref{tab:uris}.

\begin{table}[H]
\centering
\small
\begin{tabular}{@{}cll@{}}
\toprule
\textbf{Recurso} & \textbf{Verbos soportados} & \textbf{Respuesta} \\
\midrule
/api/auth/login, /register & POST & LoginResponse / 201 \\
/api/solicitudes & GET, POST & List / 201 \\
/api/solicitudes/mis-solicitudes & GET & List / 200 \\
/api/anteproyectos & GET, POST & Anteproyecto / 201 \\
/api/cronogramas & GET, POST, PUT & Cronograma / 200 \\
/api/evaluaciones & GET, POST & Evaluacion / 201 \\
/api/jurados & GET, POST & Jurado / 201 \\
/api/actas & GET, POST & Acta / 201 \\
/api/reportes/** & GET & ByteArray (PDF) / 200 \\
\bottomrule
\end{tabular}
\caption{Recursos principales de la API REST del Sistema de Pre-Sustentaciones (ruta base /api).}
\label{tab:uris}
\end{table}

\subsection{JWT: estructura y flujo de autenticación}

Un JWT (\textit{JSON Web Token}) está compuesto por tres partes separadas por puntos:
\begin{enumerate}
    \item \textbf{Header}: algoritmo de firma (HS256) y tipo de token.
    \item \textbf{Payload}: \textit{claims} (datos reclamados). En la versión actual se incluyen `sub`, `iat` y `exp`.
    Se propone ampliar a los 7 claims del RFC 7519 (`jti`, `iss`, `sub`, `aud`, `iat`, `nbf`, `exp`).
    \item \textbf{Signature}: firma HMAC-SHA256 calculada con la clave secreta (\texttt{jwt.secret}).
\end{enumerate}

El flujo implementado es el siguiente: el usuario se autentica en \texttt{POST /api/auth/login}; el servidor valida
la contraseña con \texttt{BCryptPasswordEncoder}, genera el JWT mediante \texttt{JwtTokenProvider} y devuelve
la respuesta con la cookie \texttt{jwtToken} (HttpOnly, Max-Age 24h) y el objeto \texttt{LoginResponse}.

\begin{lstlisting}[style=javaStyle, label=lst:jwt, caption={JwtTokenProvider.java — generación del token (fragmento)}]
public String generateToken(Authentication authentication) {
    UserDetails userDetails = (UserDetails) authentication.getPrincipal();
    Date now = new Date();
    Date expiryDate = new Date(now.getTime() + jwtExpiration);

    return Jwts.builder()
            .subject(userDetails.getUsername())
            .issuedAt(now)
            .expiration(expiryDate)
            .signWith(getSigningKey())
            .compact();
}
\end{lstlisting}

\subsection{OpenAPI / Swagger 3.0}

El proyecto documenta su API con \textbf{Springdoc OpenAPI 2.3.0}. La configuración (\texttt{OpenApiConfig.java})
define el título \textbf{``Sistema de Gestión de Pre-Sustentaciones UTEQ - API OpenAPI 3.0''}, la versión 0.9.0-rc
y el esquema de seguridad por token Bearer JWT. La interfaz interactiva de Swagger UI queda accesible en
\texttt{/swagger-ui/index.html}. Adicionalmente, el proyecto incluye la colección Postman exportada en
\texttt{docs/postman/PFC-Collection.json} con 22 peticiones organizadas y probadas.

\begin{lstlisting}[style=javaStyle, label=lst:oa, caption={OpenApiConfig.java — bean OpenAPI con Bearer auth (fragmento)}]
@Bean
public OpenAPI customOpenAPI() {
    final String securitySchemeName = "bearerAuth";
    return new OpenAPI()
            .info(new Info()
                    .title("Sistema de Gestión de Pre-Sustentaciones UTEQ - API OpenAPI 3.0")
                    .version("0.9.0-rc"))
            .addSecurityItem(new SecurityRequirement().addList(securitySchemeName))
            .components(new Components()
                    .addSecuritySchemes(securitySchemeName,
                            new SecurityScheme()
                                    .name(securitySchemeName)
                                    .type(SecurityScheme.Type.HTTP)
                                    .scheme("bearer")
                                    .bearerFormat("JWT")));
}
\end{lstlisting}
